% !PRINTED 258
\weberpage{258}

In dem Falle $n=1$, $p_{,}=5$ haben wir aber in $G$ nach (2) % [sic] odd subscript comma printed after p
und 3) die Substitution

\[
\begin{pmatrix} 1, & 1 \\ 2, & -2 \end{pmatrix}
\begin{pmatrix} 0, & 1 \\ -1, & 0 \end{pmatrix}
\begin{pmatrix} -2, & -1 \\ -2, & 1 \end{pmatrix}
\begin{pmatrix} 0, & 1 \\ -1, & 0 \end{pmatrix}
=
\begin{pmatrix} -2, & 0 \\ 0, & 2 \end{pmatrix},
\]
und folglich auch
\weq{12}{
\begin{pmatrix} 1, & 0 \\ -\xi, & 1 \end{pmatrix}
\begin{pmatrix} -2, & 0 \\ 0, & 2 \end{pmatrix}
\begin{pmatrix} 1, & 0 \\ \xi, & 1 \end{pmatrix}
\begin{pmatrix} 2, & 0 \\ 0, & -2 \end{pmatrix}
=
\begin{pmatrix} 1, & 0 \\ -2\xi, & 1 \end{pmatrix},
}
also, da $\xi$ und folglich auch $-2\xi$ beliebig ist, wieder die Substitutionen (10) und damit also alle Erzeugenden der Gruppe $E$,
\gesperrt{und also ist in allen diesen Fällen $G$ mit $E$ identisch und
folglich die Gruppe $E$ einfach}.

Der Fall $n=1$, $p=3$ bietet aber die zweite wirkliche
Ausnahme. In diesem Falle ist die Gruppe $E$ vom $12^{\mathrm{ten}}$ Grade
und sie hat den Normaltheiler $4^{\mathrm{ten}}$ Grades

\[
G =
\begin{pmatrix} 1, & 0 \\ 0, & 1 \end{pmatrix},
\begin{pmatrix} 0, & 1 \\ -1, & 0 \end{pmatrix},
\begin{pmatrix} -1, & 1 \\ 1, & 1 \end{pmatrix},
\begin{pmatrix} 1, & 1 \\ 1, & -1 \end{pmatrix}.
\]

Um sich davon zu überzeugen, braucht man nur die beiden
Nebengruppen:
\[
\begin{aligned}
\begin{pmatrix} 1, & 0 \\ 1, & 1 \end{pmatrix} G &=
\begin{pmatrix} 1, & 0 \\ 1, & 1 \end{pmatrix},
\begin{pmatrix} 0, & 1 \\ -1, & 1 \end{pmatrix},
\begin{pmatrix} -1, & 1 \\ 0, & -1 \end{pmatrix},
\begin{pmatrix} 1, & 1 \\ -1, & 0 \end{pmatrix} \\
\begin{pmatrix} 1, & 0 \\ -1, & 1 \end{pmatrix} G &=
\begin{pmatrix} 1, & 0 \\ -1, & 1 \end{pmatrix},
\begin{pmatrix} 0, & 1 \\ -1, & -1 \end{pmatrix},
\begin{pmatrix} -1, & 1 \\ -1, & 0 \end{pmatrix},
\begin{pmatrix} 1, & 1 \\ 0, & 1 \end{pmatrix}
\end{aligned}
\]
mit $G\begin{pmatrix} 1, & 0 \\ 1, & 1 \end{pmatrix}$, $G\begin{pmatrix} 1, & 0 \\ -1, & 1 \end{pmatrix}$
zu vergleichen.

Damit ist also allgemein bewiesen, dass, von den beiden
Fällen $n=1$, $p=2$ und $n=1$, $p=3$ abgesehen, die Gruppe $E$
einfach ist. Für die ersten Fälle erhält man für die Grade $g$
dieser einfachen Gruppen:

\[
\begin{array}{lll@{\qquad}lll}
n=1, & p=\ 5, & g=\ \ 60 & n=2, & p=2, & g=\ \ 60 \\
     & p=\ 7, & g=\ 168  &      & p=3, & g=\ 360 \\
     & p=11,  & g=\ 660  &      & p=5, & g=7800 \\
     & p=13,  & g=1092   & n=3, & p=2, & g=\ 504 \\
     &        &          &      & p=3, & g=9828 \\
     &        &          & n=4, & p=2, & g=4080^{1)}.
\end{array}
\]

\rule{4cm}{0.4pt}

$^{1)}$ Siehe Bulletin of the New York math. Society. October 1893. In
dem Berichte über den mathematischen Congress in Chicago finden sich
zwei Notizen über diese Gruppen von \gesperrt{Cole} und \gesperrt{Moore}.
